\section{从实数到浮点数}
\label{sec:08-Numerical-Analysis/01-Floating-Point-and-Error/01}

打开任何一个编程环境，敲一行最简单的算术：
\begin{verbatim}
0.1 + 0.2
\end{verbatim}
你预期看到 \(0.3\)。但你看到的很可能是 \(0.30000000000000004\)。小数点末位多出一个 4，不多不少，每次运行都一样。

这不是 bug。你的 CPU 没有算错，编译器也没有偷懒。它只是诚实地告诉了你一件事：程序里从来没有出现过真正的 \(0.1\)，也从来没有出现过真正的 \(0.2\)。

\subsection{二进制展开的困境}

十进制写 \(0.1\) 只需要一位小数。但二进制里，\(0.1\) 的展开是无限循环的：
\[
(0.1)_{10} = (0.00011001100110011\ldots)_2
\]
把十进制的「十分之一」放进二进制的世界，就像把 \(1/3\) 写成十进制小数——你永远写不完，只能在某一位截断。

所以程序里存的三个数，从一开始就只是三个近似值：
\begin{itemize}
\tightlist
\item 名叫 \verb|0.1| 的那个二进制数，不等于数学上的十分之一；
\item 名叫 \verb|0.2| 的那个二进制数，也不等于数学上的十分之二；
\item 名叫 \verb|0.3| 的那个二进制数，同样不等于数学上的十分之三。
\end{itemize}
把前两个近似值加起来，再舍入到最近的浮点数，碰巧不等于第三个近似值——于是多出一个末位 4。

这件事不该被概括为「计算机算错了小学加法」。它更应该被理解成：在有限精度系统中，「相等」的定义变了。

\subsection{什么时候该在乎}

同样是写 \(0.1\)，不同场景对精度的需求完全不同。

财务系统里，0.1 很可能代表一角钱。十元钱减去三笔一角钱，账户必须精确归零。让二进制近似误差来决定「归零了没」，迟早会在审计报表上炸出一个本该不存在的余额。常见的做法是：金额按最小货币单位存成整数（分），或者使用十进制算术类型，明确指定舍入规则。

传感器读数里的 0.1 是另一回事。一个温度计读数 37.1 度，它本身就带着±0.05 度的测量误差。你用 \verb|abs(t1 - t2) < 1e-6| 的方式比较两个测量值，反而比要求逐位相等更合理。

判断之前，先问自己：这个 0.1 是一个精确量，还是一个近似量？两种答案指向完全不同的数值表示策略。

\subsection{浮点数可以看成有限位的科学计数法}

刚才的困境来自「位不够」。浮点数系统的回应是：用科学计数法的方式，把有限的位分配到需要的地方。一个规格化二进制浮点数写作
\[
x = (-1)^s \times (1.f)_2 \times 2^e,
\]
其中 \(s\) 控制正负，\(f\) 是有效数字的二进制小数部分，\(e\) 是指数。

你可以把它想象成：有效位是固定长度的尺子，指数告诉你这把尺子放在哪个数量级上。同一个长度尺，放在 1 附近能量出 \(2^{-52}\) 的间距；放在 \(2^{53}\) 附近，相邻刻度之间的距离已经是 2。

这就是浮点系统的核心权衡：用指数的范围换取覆盖的广度，代价是精度随尺度变化。

\subsection{精度近似恒定的是相对量，不是绝对量}

在区间 \([2^e, 2^{e+1})\) 内，相邻规格化数的间距是固定的。一旦进入下一个二进制数量级，间距立刻翻倍。

以 IEEE 754 binary64（常见的「双精度」）为例：有 53 个二进制有效位。在 1 附近，相邻浮点数的间距是
\[
2^{-52} \approx 2.22 \times 10^{-16}.
\]
这个数字常被翻译成「双精度有大约 16 位十进制有效数字」。这句话没错，但很容易被误解。

我们来做一个实验。试试把 \(2^{53} + 1\) 存进 binary64。结果你会发现：
\[
2^{53} + 1 = 2^{53}.
\]
两个相差 1 的整数，在计算机里变成了同一个数。不是因为 1 太小——1 在它自己的尺度上很大——而是因为在 \(2^{53}\) 这个数量级上，相邻浮点数的间距已经是 2 了。加上 1，离下一个可表示数的距离还差一半，舍入到最近就回到原地。

换一个方向看：在此之下的很多整数是可以精确表示的。\(2^{53} - 1\) 可以，\(2^{53} - 1000\) 也可以。但从 \(2^{53}\) 往上，整数开始出现「缺口」——不是所有整数都能找到对应的浮点数。到 \(2^{54}\)，相邻间距变成 4；继续往上，缺口越来越大。

浮点系统提供的是近似恒定的相对精度，而不是恒定的绝对精度。把 10 亿算错 1 块钱相对误差十万亿分之一；把 0.0000001 算成 0 则可能完全改变结论，尽管绝对误差才千万分之一。

\subsection{machine epsilon 与 unit roundoff}

谈到浮点精度时，你会在不同书里看到两个相近但不同的数字。混用它们会导致误差分析差一个因子 2。

\begin{itemize}
\tightlist
\item \(\varepsilon_{\mathrm{mach}}\)：从 1 到下一个更大浮点数的间距。在 binary64 中，这个值是 \(2^{-52}\)。
\item unit roundoff \(u\)：舍入到最近时的最大相对舍入误差。在 binary64 中，最坏情况相对误差是半个间距除以 1，即 \(2^{-53} = \varepsilon_{\mathrm{mach}} / 2\)。
\end{itemize}

理解这两个数字不需要记住二进制字段排布。你只需要知道：\(u\) 是你对一次基本运算能期待的相对精度的上限，前提是结果没有溢出也没有下溢。

有了 unit roundoff，我们可以写出浮点运算的标准模型。对于舍入到最近的情况：
\[
\operatorname{fl}(x) = x(1 + \delta), \qquad |\delta| \le u.
\]
对一次基本运算 \(\circ \in \{+, -, \times, \div\}\)：
\[
\operatorname{fl}(a \circ b) = (a \circ b)(1 + \delta), \qquad |\delta| \le u,
\]
假设精确结果在正常范围内。

这个模型的直观含义是：每次运算就像精确计算后再做一次微小的相对扰动，扰动不超过 \(u\)。它是一个分析工具，不是全实轴上的无条件定理——特殊值、溢出、下溢都不在它的覆盖范围内。

\subsection{零附近发生了什么}

如果只在每个二进制数量级里均匀排布规格化数，零和最小的规格化数之间会留出一个巨大的空档。任何比最小规格化数还小的量，会直接下溢成零。

IEEE 754 用次正规数（subnormal numbers）填补了这片区域。它们均匀地排在零和最小规格化数之间，让间距逐渐延伸到零。这叫 gradual underflow——缓慢下溢。代价是：次正规数区域的相对精度不再保持；一个次正规数只有较少的有效位，因为高位的前导 1 被移走了。

除了次正规数，标准还定义了几种特殊值，各有各的用途：

\begin{itemize}
\tightlist
\item 正零与负零：保留极限方向的符号信息。\(1/\infty\) 是 \(+0\)，\(1/(-\infty)\) 是 \(-0\)，尽管它们比较相等。
\item 正负无穷：来自非零有限数除以零，或正常运算溢出。它们参与后续运算时有明确规则，比如 \(1/\infty = 0\)。
\item NaN（Not a Number）：表示未定义或无效结果，如 \(0/0\)、\(\infty - \infty\)、\(\sqrt{-1}\)。NaN 不跟任何值比较相等，包括它自身。这个设计故意让 \verb|x == NaN| 永远为 false；检测 NaN 需要用专门的函数。
\item 多种舍入模式：默认是舍入到最近且 ties-to-even（正好在中间时优先选末位为偶数的那个），还有向零、向正无穷、向负无穷舍入等选项。
\end{itemize}

无穷和 NaN 让异常得以传播而不是悄无声息地崩溃。但它们不能替代输入检查——让 NaN 在计算图中流动数百步后再追溯来源，比一开始就检查困难得多。

\subsection{浮点运算不是实数代数}

我们在学校学的算式操作里有一大串「放之四海皆准」的规律：加法结合律、乘法分配律、同类项合并。在浮点世界，这些规律通通变成了「有时成立」。

结合律是最先倒下的：
\[
(a + b) + c \ne a + (b + c).
\]
原因很简单：每次加法之后都会舍入。如果 \(a\) 非常大而 \(b\) 非常小，\(a+b\) 可能舍入回 \(a\)。再加 \(c\) 的结果，跟先把 \(b\) 和 \(c\) 加起来再加到 \(a\)，中间丢失的信息不同，最终答案也就不同。

改变求和顺序会改变哪些低位先被丢弃。乘法对加法的分配律也可能失效——两个理论上相等的表达式，算出来的浮点结果偏偏有出入。

这意味着两件事。坏消息是：你不能拿起纸笔推导出来的代数等价式就往代码里抄，觉得「数学上一样所以结果应该一样」。好消息是：代数改写本身可以成为数值算法设计的一部分——选择哪个等价表达式，就是选择哪条计算路径能让误差在可控范围内。

\subsection{精确可表示与精确可计算是两件事}

\(1/8 = 0.001_2\) 可以精确表示为二进制浮点数。\(1/10\) 不可以。

但即使输入都可以精确表示，精确答案也可以精确表示，中间计算仍然可能溢出或舍入。反过来，一连串近似中间值也可能在最后一刻撞进正确的浮点数，就像盲人走迷宫恰好走到出口。

判断一次计算是否可靠，不能只问「输入能精确表示吗」。你得追踪五件事：

\begin{itemize}
\tightlist
\item 所有中间量的数量级——有没有某个步骤让两个极大的数相乘并溢出？
\item 运算的顺序——结合律不成立，哪个先算哪个后算？
\item 消去风险——哪里有相近大数相减？
\item 问题的条件性——输入末位的一点变动，答数会变多少？
\item 算法的稳定性——算法的每一步舍入积累，是不是跟问题本身的条件性在一个量级？
\end{itemize}

浮点数不是「带一点误差的实数」。它是一套有限刻度的标尺，疏密随位置而变化。你只要承认自己只能在离散点上操作，许多奇怪的结果就不再奇怪——它们只是在告诉你，你已经走到了当前精度能分辨的边界。